%O código seguinte será incluído no pdf gerado
%Mais info: http://www.tug.org/applications/hyperref/manual.html
\hypersetup {
pdfmenubar=true,	%Mostrar barra menus
pdfpagemode=UseNone, %Nenhum modo de páginas
pdfpagelayout=SinglePage, %Abertura em modo de página simples
pdffitwindow=true, %Adaptar página à janela
%
bookmarksopen=true,
pdftitle=\thesistitlestr, %Título
pdfauthor=\thesisauthorsstr, %Autor
pdfsubject=\thesistypestr, %Assunto
pdfhighlight=/O, %Efeito ao clicar link
linkcolor=linkcol, %Cor das ligações internas do documento
citecolor=citecol, %Cor das referências à bibliografia no texto
urlcolor=linkcol   %Cor das hiperligações
}

%%O comando seguinte insere o espaçamento de 1.5 linhas
\onehalfspacing

%%Página de rosto
\pagenumbering{roman}

%Início-----------------------

\begin{titlepage}
    \begin{center}
        \begin{flushright}
            \vspace*{0.03cm}\includegraphics[height=3cm]{logo_ubi_vprincipal}\\
            %\rostoubi UNIVERSIDADE DA BEIRA INTERIOR\\
            %\rostofac \facultynamestr\\
            \vspace{5.6cm}%7.6/6.6 - 1

            \rostotitulo \textbf{\thesistitlestr} \\
            \ifdefined\thesissubtitlestr
                \rostosubtit \textbf{\thesissubtitlestr}\\
            \fi
            \vspace{1.6cm}%1.8
            \rostonomes \textbf{\thesisauthorsstr} \\
            \vspace{1.4cm}%1.4
            % \rostooutros Tese para obtenção do Grau de Doutor em\\
            \rostooutros \thesistypestr \\
            \rostonomes \textbf{\thesiscoursestr}\\
            \rostooutros (\studiescyclenumberstr \ \studiescycle)\\
            \vspace{3.3cm}%3.2
            %\rostooutros \thesissupervisorsstr

            %\rostonomesdeorientor Orientador: Name of your supervisor\\
            \thesissupervisorsstr
            %Co-orientador: Prof. Doutor Nome Completo\\
            \vspace{1.4cm}

            %\rostooutros
            \rostooutros \textbf{\thesislocalanddatestr}
        \end{flushright}
    \end{center}
\end{titlepage}

%Fim------------------
%\dominitoc

%%Numeração das páginas
\pagestyle{fancy}


%%O comando a seguir gera uma página após a de rosto com cabeçalho e rodapé
\cleardoublepage

%%O comando a seguir permite que as costas da página de rosto não inclua cabeçalho mas rodapé (escolher entre este e outro)
%\newpage\mbox{}\thispagestyle{plain}\fancyhead{}